%!TEX TS-program = pdflatex
%% Modèle d'article ACM (acmart, ex. ISSAC) (Olivier Bournez). Compilation : make.
%% acmart est trop rigide pour macros.tex : ce modèle est autonome, avec un jeu minimal
%% de macros. Soumission anonyme : \documentclass[sigconf,screen,review,anonymous]{acmart}
\documentclass[sigconf,screen]{acmart}

%% Bascule finale : \FINAL{...} n'apparaît qu'en version finale, \PASFINAL{...} qu'avant.
\providecommand\FINAL[1]{}
\providecommand\PASFINAL[1]{#1}
\FINAL{\renewcommand\PASFINAL[1]{}}
\newcommand\SHORTER[1]{}
\providecommand\BIBFILES{@@reference-biblio}

%% Métadonnées ACM (fournies par la conférence à la camera-ready)
\setcopyright{cc}
\setcctype{by}
\copyrightyear{2026}
\acmYear{2026}
\acmDOI{}
\acmISBN{}
\acmConference[ACRONYME '26]{Nom complet de la conférence}{Dates}{Lieu}
\acmBooktitle{Nom complet de la conférence (ACRONYME '26), Dates, Lieu}

\AtBeginDocument{\providecommand\BibTeX{{Bib\TeX}}}
\let\Bbbk\relax
\usepackage{todonotes}
\usepackage[capitalise]{cleveref}
\usepackage{xfrac}
\usepackage{amsmath}

%% Macros minimales (celles de macros.tex qui servent vraiment)
\newcommand{\tu}[1]{\mathbf{#1}}
\newcommand{\cp}[1]{\operatorname{#1}}
\newcommand\R{\mathbb{R}}
\newcommand\N{\mathbb{N}}
\newcommand{\norm}[1]{\left\lVert#1\right\rVert}
\newcommand\Ptime{\cp{PTIME}}
\newcommand\Pspace{\cp{PSPACE}}
\newcommand\todoinline[1]{\FINAL{\todo[inline]{#1}}}

%% Environnements (acmart charge amsthm)
\newtheorem{proposition}{Proposition}
\newtheorem{theorem}{Theorem}
\newtheorem{lemma}{Lemma}
\newtheorem{corollary}{Corollary}
\theoremstyle{definition}
\newtheorem{definition}{Definition}
\newtheorem{example}{Example}
\theoremstyle{remark}
\newtheorem{remark}{Remark}

%% apxproof avec acmart : possible, décommenter (theoremrep + proofsketch + proof).
%\providecommand\apxparameter{appendix=append}
%\usepackage[\apxparameter]{apxproof}
%\newtheoremrep{theorem}{Theorem}
%\newtheoremrep{lemma}[theorem]{Lemma}

\begin{document}

\title{Titre de l'article}
\author{Olivier Bournez}
\affiliation{%
  \institution{LIX, École polytechnique, CNRS}
  \city{Palaiseau}
  \country{France}}
\email{bournez@lix.polytechnique.fr}
%% \author{Prénom Nom} \affiliation{...} \email{...}
\renewcommand{\shortauthors}{Bournez}

\begin{abstract}
Quatre phrases : la question, pourquoi elle compte, le résultat principal, la technique.
\end{abstract}

\begin{CCSXML}
<ccs2012>
<concept>
<concept_id>10003752.10003753</concept_id>
<concept_desc>Theory of computation~Models of computation</concept_desc>
<concept_significance>500</concept_significance>
</concept>
</ccs2012>
\end{CCSXML}
\ccsdesc[500]{Theory of computation~Models of computation}
\keywords{mot-clé 1, mot-clé 2, mot-clé 3}

\maketitle

\section{Introduction}
\label{sec:intro}

Contexte, question, résultat principal (\cref{thm:main}), travaux antérieurs~\cite{JournalACM2017}.

\section{Preliminaries}
\label{sec:prelim}

\begin{definition}[Notion]
\label{def:notion}
Une définition.
\end{definition}

\section{Main results}
\label{sec:results}

\begin{theorem}[Main theorem]
\label{thm:main}
Énoncé du théorème principal.
\end{theorem}
\begin{proof}
Preuve.
\end{proof}

\section{Conclusion and perspectives}
\label{sec:conclusion}

Ce qui est fait, ce qui reste ouvert.

\bibliographystyle{ACM-Reference-Format}
\bibliography{\BIBFILES}
\end{document}
